\documentclass[11pt]{article}
\usepackage{amsmath,amssymb,amsthm}
\usepackage[T1]{fontenc}
\usepackage{lmodern}
\usepackage{hyperref}
\usepackage{microtype}
\setlength{\textwidth}{6.15in}
\setlength{\oddsidemargin}{0.17in}
\setlength{\evensidemargin}{0.17in}
\setlength{\textheight}{8.8in}
\setlength{\topmargin}{-0.28in}
\newtheorem{theorem}{Theorem}
\newtheorem{corollary}[theorem]{Corollary}
\newtheorem{remark}[theorem]{Remark}
\title{A Sharp Weight-Ratio Bound for the Anisotropic\\
Navier--Stokes Terminal-Core Expansion}
\author{Tony Newton}
\date{}
\begin{document}
\maketitle
\begin{abstract}
The convergence argument for the terminal-core coefficient hierarchy in
\emph{Exact Local Unforced Terminal-Core Dynamics in Three-Dimensional
Navier--Stokes} uses a radial regular-singular inverse to compensate for
two derivatives in an axial similarity parameter.  The original analytic
estimate bounded the relevant weight ratio by $16R/\rho^2$.  Here we obtain
the sharp coefficient bound $8R/\rho^2$, uniformly in both derivative
indices and for each of the two radial inverses.  We derive the associated
operator inequality on the explicitly defined weighted coefficient space
and distinguish this estimate from the coupled higher-order inverse
needed for convergence of the complete Navier--Stokes series.
\end{abstract}

\section{The regular-singular estimate}
The local unforced construction in \cite{newton2026} introduces formal
coefficients of a terminal similarity expansion.  Its higher-order
equations contain two derivatives in the axial similarity coordinate
$\eta$.  A useful bound must therefore account for the \emph{same}
analytic norm at every expansion order, rather than successively
reducing a strip of analyticity.

For a formal radial series
$F(X,\eta)=\sum_{\alpha\geq0}F_\alpha(\eta)X^\alpha$, and fixed $R,\rho>0$,
write
\[
 w_{\alpha,\beta}
 =\frac{R^{-\alpha}\rho^{-\beta}\beta!
   \binom{2\alpha+\beta}{\beta}}
 {(\alpha+1)^2(\beta+1)^2},\qquad
 \|F\|_{R,\rho}
 =\sup_{\alpha,\beta\geq0}
 \frac{\lvert\partial_\eta^\beta F_\alpha\rvert}
 {w_{\alpha,\beta}}.
\]
When the coefficient derivatives exist and this supremum is finite,
we consider the axis-regular inverse of
\[
 XG_{XX}+\nu G_X=F,\qquad G(0,\eta)=0,
 \qquad \nu\in\{1,2\}.
\]
Coefficientwise it is
\[
 (J_\nu F)_{\alpha+1}
  =\frac{F_\alpha}{(\alpha+1)(\alpha+\nu)}.
\]
The factor with which $J_\nu\partial_\eta^2$ changes the weighted
coefficient seminorm is exactly
\begin{equation}\label{eq:ratio}
 \frac{w_{\alpha,\beta+2}}
 {(\alpha+1)(\alpha+\nu)\,w_{\alpha+1,\beta}}
 =\frac{R}{\rho^2}
 \frac{2(\alpha+2)^2(2\alpha+1)(\beta+1)^2}
 {(\alpha+1)^2(\alpha+\nu)(\beta+3)^2}.
\end{equation}
The estimate in the original paper takes the dimensionless last
fraction to be at most $16$.  The optimal uniform bound is $8$.

\begin{theorem}[Sharp uniform two-derivative ratio]\label{thm:sharp}
For all real $a,b\geq0$ and $\nu\geq1$,
\begin{equation}\label{eq:sharp}
 \frac{2(a+2)^2(2a+1)(b+1)^2}
 {(a+1)^2(a+\nu)(b+3)^2}\leq8.
\end{equation}
No smaller constant works uniformly for nonnegative integral indices
$a,b$ and $\nu\in\{1,2\}$.
\end{theorem}

\begin{proof}
The identity
\[
 4(a+1)^3-(a+2)^2(2a+1)=a^2(2a+3)
\]
has nonnegative right-hand side for $a\geq0$.  Since
$a+\nu\geq a+1$ and $(b+1)^2\leq(b+3)^2$, the fraction in
\eqref{eq:sharp} is at most
\[
 \frac{2\cdot4(a+1)^3}{(a+1)^3}
 \frac{(b+1)^2}{(b+3)^2}
 \leq8.
\]
For $a=0$ and $\nu=1$ the fraction is precisely
$8(b+1)^2/(b+3)^2$, which tends to $8$ as $b\to\infty$
through the nonnegative integers.
\end{proof}

\begin{corollary}[Two-derivative inverse]\label{cor:inverse}
In the weighted coefficient space above,
\[
 \|J_\nu\partial_\eta^2F\|_{R,\rho}
 \leq \frac{8R}{\rho^2}\|F\|_{R,\rho},
 \qquad \nu\in\{1,2\},
\]
whenever the coefficientwise operation is defined.
\end{corollary}

\begin{proof}
The $\beta$th derivative of the $(\alpha+1)$st coefficient of
$J_\nu\partial_\eta^2F$ equals
$\partial_\eta^{\beta+2}F_\alpha/
 ((\alpha+1)(\alpha+\nu))$. Divide by $w_{\alpha+1,\beta}$,
apply \eqref{eq:ratio} and Theorem~\ref{thm:sharp}, and take
the supremum over both indices.  The constant radial
coefficient of the image vanishes.
\end{proof}

\section{Relation to the full hierarchy}
Corollary~\ref{cor:inverse} improves the constant in
Lemma~4.2 and Appendix~B.2 of \cite{newton2026}.  The improvement
does not change the analytic space, the two-derivative allowance,
or the radial valuation assigned to each expansion order.

There are further hypotheses in the passage from the scalar inverse
to the Navier--Stokes coefficient hierarchy.  In particular, the
current-order feedback $K_n$ is a coupled operator involving the
radial and axial velocity coefficients, the similarity shifts,
and pressure integration.  The estimate
\[
 \|K_n\|_{I_n\times I_n\to I_n\times I_n}\leq \frac{C}{n}
\]
depends on controlling all four current-order blocks on a common
analytic space.  Neither Theorem~\ref{thm:sharp} nor
Corollary~\ref{cor:inverse} supplies those four block estimates
by itself.  The exponential majorant and exact unforced physical
reconstruction consequently remain separate steps in a formal
verification of the main convergence theorem.

\section{Formal verification}
The inequality \eqref{eq:sharp}, stated for all real
$a,b\geq0$ and $\nu\geq1$, is formalized in Lean~4 as
\texttt{TerminalCore2026.two\_derivative\_ratio\_le\_eight}.
This is stronger than checking finitely many triples of indices:
its proof relies on the polynomial identity used above and
ordered-real arithmetic.  The accompanying Lean modules also
check coefficientwise radial inversion, scaling-exponent identities,
and a conditional first-zero obstruction for the axial derivative.
The source, pinned toolchain, and GitHub Actions jobs appear at
\begin{center}
\small
\url{https://github.com/tonynewton79-web/unforced-navier-stokes-lean/pull/1}.
\end{center}
The Lean theorem directly certifies the universal
\emph{ratio inequality}.  A separate Lean implementation of the
weighted coefficient space and the operators remains necessary
to kernel-check Corollary~\ref{cor:inverse} as a function-space result.

\begin{thebibliography}{9}
\bibitem{newton2026}
T. Newton,
\emph{Exact Local Unforced Terminal-Core Dynamics in
Three-Dimensional Navier--Stokes: Anisotropic Concentration and Swirl
Reversals}, Zenodo (2026),
\url{https://doi.org/10.5281/zenodo.23474915}.
\end{thebibliography}
\end{document}
